\BourbakiExercises{习题}

\begin{enumerate}
\def\labelenumi{\arabic{enumi})}
\tightlist
\item
  \begin{enumerate}
  \def\labelenumii{\alph{enumii})}
  \tightlist
  \item
    设 M 是一个 A-模，x 是 M 中零化子为有限族极大理想之交的元素。证明模 Ax 是半单的（把它嵌入到单模的一个有限乘积中）。
  \end{enumerate}
\end{enumerate}

\begin{enumerate}
\def\labelenumi{\alph{enumi})}
\setcounter{enumi}{1}
\tightlist
\item
  由此推出：一个模 M 是半单的当且仅当每个非零元素的零化子都是有限族极大理想之交。
\end{enumerate}

\begin{enumerate}
\def\labelenumi{\arabic{enumi})}
\setcounter{enumi}{1}
\tightlist
\item
  设 M 是一个模。证明下列性质等价：
\end{enumerate}

\(\alpha)\) M 是半单的且长度有限。

\(\beta)\) M 是诺特的，且 M 的每个极大子模都有补。

\(\gamma)\) M 是阿廷的，且 M 的每个单子模都有补。

\begin{enumerate}
\def\labelenumi{\arabic{enumi})}
\setcounter{enumi}{2}
\item
  设 \(\mathrm{K}\) 是一个体，I 是无限集，A 是环 \(\mathrm{K}^{\mathrm{I}}\)。证明 A-模 \(\mathrm{A}_s\) 不是半单的，尽管每个单生成理想都有补（确定 \(\mathrm{A}_s\) 的所有单子模）。
\item
  设 M 是一个模，它是子模族 \((\mathrm{M}_{\alpha})_{\alpha\in\mathrm{I}}\) 的直和。证明子模 \(M_{\alpha}\) 在 M 的所有自同构下稳定当且仅当每个同态 \(M_{\alpha}\to M_{\beta}\)（其中 \(\alpha\neq\beta\)）都为零。
\item
  设 M 是一个有限长度的模，\(M = \oplus_{i \in I} M_i\) 是 M 分解为不可分解子模直和的一个分解。设 \(J \subset I\) 是 i 与 j 之间关系“模 \(M_i\) 与 \(M_j\) 同构”的一个等价类；令 \(M_J = \oplus_{i \in J} M_i\)。举一个例子表明 \(M_J\) 不一定在 M 的每个自同构下稳定（取 M 为有限长度的 Z-模并利用习题 4）。
\item
  若模 M 的每个非零商都含一个单子模，则称 M 为半阿廷模。每个阿廷模以及体上的每个向量空间都是半阿廷的。
\end{enumerate}

\begin{enumerate}
\def\labelenumi{\alph{enumi})}
\item
  设 \(0 \rightarrow M' \rightarrow M \rightarrow M'' \rightarrow 0\) 是一个正合列。则 M 是半阿廷的当且仅当 \(M'\) 与 \(M''\) 亦然（若 M 是半阿廷的，考虑由满足 \(N \cap M' = 0\) 的 M 的子模 N 组成的集合中的一个极大元素）。
\item
  一族半阿廷模的直和是半阿廷模。 c) 设 K 是一个体，A 是乘积环 \(K^{N}\)。证明 A-模 \(A_{s}\) 不是半阿廷的，尽管它是单 A-模的一个乘积（若 s 是 \(A_{s}\) 的底座，证明 A-模 \(A_{s}/s\) 不含任何单模）。
\item
  设 A 是左诺特环，M 是半阿廷 A-模。证明 M 是它的阿廷子模之和（设 S 是 M 的阿廷子模之和；注意，若 N 是 M 的在 M/S 中有单像的有限生成子模，则 \(N \cap S\) 是阿廷的）。特别地，若 M 是有限生成的，则它是阿廷的。
\item
  设 D 是一个体，V 是无限维 D-向量空间，A 是由 D 的中心与有限秩自同态生成的 \(\operatorname{End}_{\mathrm{D}}(\mathrm{V})\) 的子环。证明 \(A_{s}\) 是半阿廷的（利用 a) 与 b)），但不是它的阿廷子模之和。
\end{enumerate}

\begin{enumerate}
\def\labelenumi{\arabic{enumi})}
\setcounter{enumi}{6}
\tightlist
\item
  设 A 是一个环，M 是一个 A-模。若 M 的子模 N 对 M 的每个非零子模 P 都满足 \(N \cap P \neq 0\)，则称 N 为本质子模。
\end{enumerate}

\begin{enumerate}
\def\labelenumi{\alph{enumi})}
\item
  证明 M 的每个子模 N 都是某个本质子模的直因子（考虑一个在与 N 之交为零的诸子模中极大的子模 P）。
\item
  证明 M 的底座是 M 的诸本质子模之交（由 a) 推出这个交是半单的）。特别地，M 是半单的当且仅当它不含除自身以外的本质子模。
\end{enumerate}

¶ 8) 设 A 是左诺特环。

\begin{enumerate}
\def\labelenumi{\alph{enumi})}
\item
  设 x 是 A 的非零元素。证明 x 的左零化子 Ann x 不是本质的（在不具有此性质的元素中，取一个使 Ann x 极大的 x，并考虑 \(Ax \cap Ann x\)）。
\item
  设 x 是 A 的右可约元。证明理想 Ax 是本质的（若 a 是 A 的理想且满足 \(a \cap Ax = 0\)，考虑理想 \(a + ax + ax^{2} + \ldots\)）。
\item
  由 a) 与 b) 推出，A 的每个右可约元都是可约元。
\item
  设 a 是 A 的一个含非幂零元素的本质左理想。证明 a 含一个可约元（设 \((x_{1},\ldots,x_{n})\) 是 a 中满足 Ann \(x_{i}^{2}=\) Ann \(x_{i}\) 且 \(x_{i+1}\in\) Ann \(x_{i}\cap\cdots\cap\) Ann \(x_{1}\)（对每个 \(i\geqslant1\)）的非零元素序列；注意 a 含诸 \(Ax_{i}\) 的直和。由此推出存在这样的序列使得 \(\bigcap_{i}\) Ann \(x_{i}=0\)；证明元素 \(\sum x_{i}^{2}\) 这时是右可约的，并利用 c)）。
\end{enumerate}

